\section{この結果で言えないこと}
\label{sec:limits}

\begin{enumerate}
  \item \textbf{確率は過去 10 年の割合である}。将来の冬が同じ割合になる保証はない。
    年による違い（図~\ref{fig:byyear}）が \SI{\pm 10}{\percent} ほどあるので、
    1 冬で経験する割合はこの範囲で上下する。
  \item \textbf{予報が当たる前提で数えている}。実際の判断は前日から当日朝の予報で行う。
    予報の落下点がずれて、目標に入るはずが山に落ちる日や、その逆の日が出る。
    予報の誤差はこの計算に入っていない。目標の縁に近い日ほど、予報と実際で成否が入れ替わりやすい。
  \item \textbf{地上気象は \ang{0.25}（約 \SI{25}{\kilo\metre}）格子の値である}。
    島や海岸では、格子に海が混ざるかどうかで地上の風が大きく変わる（\ref{sec:amedas} 節）。
    本文では最寄りの陸の格子を使ったが、実際の打ち上げ場所の風は地形や建物で
    さらに変わる。現地の風の記録で確かめる必要がある。
  \item \textbf{9 時の 1 回だけを見ている}。朝の風が強くても昼に弱まる日は、実際には
    時刻をずらして上げられる。そのぶん実施できる確率は本文の値より高い可能性がある。
    逆に、日没までに回収を終える時間の制約は入れていない。
  \item \textbf{上空の風は JRA-3Q の \ang{0.375} 格子・6 時間ごとの解析値である}。
    飛行中の風の時間変化（約 2 時間）は入れていない。
  \item \textbf{瀬戸内海への着水の扱い}。瀬戸内海に落ちた機体を回収できるか、
    船舶の多い海域に落としてよいかは、この計算の外の問題である。
    上位の候補地では成功の 7--9 割が瀬戸内海への着水なので、これを認めるかどうかで結論が変わる。
  \item \textbf{平野の定義は地形だけで決めている}。水田・畑・森林・河川・道路の区別はしていない。
    冬の水田は回収しやすいが、ため池や河川に落ちることもある。
  \item \textbf{航空法などの手続き}は考えていない。空港の周辺や航空路の下では
    飛行の許可が必要で、候補地によっては条件が付く。
\end{enumerate}

\section*{データの出典}
\begin{itemize}
  \item 上空の風: JRA-3Q（気象庁）。NSF NCAR GDEX（d640000）から取得。
  \item 積雪・地上気象: ERA5（Copernicus Climate Change Service）。NSF NCAR GDEX（d633000）から取得。
    Copernicus Climate Change Service の情報を加工したものを含む。
  \item 地上の観測: 気象庁 アメダス（過去の気象データ）。
  \item 地形: NOAA NCEI DEM mosaic。
  \item 市街地: 国土交通省 国土数値情報（人口集中地区データ、2020 年）。
  \item 広場: \textcopyright{} OpenStreetMap contributors（ODbL）。
\end{itemize}
JRA-3Q から作った図と数値は、元データの条件（CC BY-NC-SA 4.0、非商用）に従う。
